\documentclass[11pt]{article}
\usepackage[utf8]{inputenc}
\usepackage{amsmath,amssymb}
\usepackage[margin=1in]{geometry}
\usepackage{hyperref}
\emergencystretch=3em

\title{The fluctuation ladder: lowering $2S'_\lambda-\beta a S_\lambda=(2\lambda-1)S_{\lambda-1/2}$}
\author{Claude, with Daniel Murfet}
\date{2026-09-06}

\begin{document}
\maketitle
\sloppy

\begin{abstract}
Unit 5 of the grammar-paper \S4 autoformalisation: the lowering identity of the fluctuation-function
ladder algebra (\texttt{lem:oscillator\_algebra}), $2S'_\lambda(a)-\beta a\,S_\lambda(a)
=(2\lambda-1)S_{\lambda-1/2}(a)$ for $\lambda>1/2$. It is the $\beta$-power-free form of the paper's
$b\,S_\lambda=(2\lambda-1)\beta^{-1/2}S_{\lambda-1/2}$, and falls out of property (i) and the
recurrence in a few lines. Zero \texttt{sorry}/\texttt{axiom}; proven first try.
\end{abstract}

\section{Setting}
The ladder algebra of \S4 has a raising operator $b^\dagger=\beta^{-1/2}\partial_a$ and a lowering
operator $b=2\beta^{-1/2}\partial_a-\beta^{1/2}a$. The raising action $b^\dagger S_\lambda=\beta^{1/2}
S_{\lambda+1/2}$ is exactly the derivative ladder already formalised (\texttt{deriv\_fluctuation},
unit~2). The lowering action is the one new identity, and it is what makes the pair a genuine
oscillator algebra. Multiplying the paper's statement through by $\beta^{1/2}$ removes every
fractional power of $\beta$, giving the clean form above.

\section{The thing formalised}
{\small
\begin{verbatim}
theorem fluctuation_lowering (β lam : ℝ) (hβ : 0 < β) (hlam : 1 / 2 < lam) (a : ℝ) :
    2 * deriv (fun a => fluctuation β lam a) a - β * a * fluctuation β lam a
      = (2 * lam - 1) * fluctuation β (lam - 1 / 2) a
\end{verbatim}
}

\section{Proof strategy}
By the derivative ladder, $S'_\lambda=\beta S_{\lambda+1/2}$, so the left side is
$2\beta S_{\lambda+1/2}-\beta a S_\lambda$. The recurrence at $\lambda-\tfrac12$ (valid since
$\lambda>\tfrac12$) reads $S_{\lambda+1/2}=\tfrac{a}{2}S_\lambda+\tfrac{\lambda-1/2}{\beta}
S_{\lambda-1/2}$; substituting cancels the $\beta a S_\lambda$ term and leaves
$2(\lambda-\tfrac12)S_{\lambda-1/2}=(2\lambda-1)S_{\lambda-1/2}$. In Lean: rewrite the derivative,
rewrite the recurrence (with the index normalisations $(\lambda-\tfrac12)+1=\lambda+\tfrac12$ and
$(\lambda-\tfrac12)+\tfrac12=\lambda$), then \texttt{field\_simp; ring}.

\section{Roadblocks and resolutions}
None --- first try. The only care was normalising the recurrence's shifted indices to match
$S_{\lambda+1/2}$ and $S_\lambda$ before substituting.

\section{What was Mathlib, what was new}
No Mathlib beyond \texttt{field\_simp}/\texttt{ring}; the content is purely the two prior grammar
results (derivative ladder + recurrence) combined.

\section{Lessons}
When a paper states an operator identity with fractional powers of a parameter, clearing them by
scaling the whole equation (here $\times\beta^{1/2}$) gives a Lean statement free of
\texttt{rpow}-on-$\beta$ bookkeeping while remaining faithful --- the operator content is unchanged.

\section{Follow-ups}
Remaining \S4: \texttt{prop:fluctuation\_weber} (the ODE~(iii) $\to$ Weber's equation substitution;
formalisable as an ODE transform, no special-function library needed) and the commutator
$[b,b^\dagger]=1$ (needs an operator framework). \texttt{cor:fluctuation\_closed\_form} and
\texttt{lem:qho} rest on parabolic-cylinder functions $D_\nu$, which Mathlib does not have --- deferred.
\end{document}
